\section{Ablation Scope}
\label{sec:ablation}

The B1, B4, B6, and B7 ablation results are used as diagnostic evidence for interpreting the effect of ArcFace, BNNeck, and light SupCon supervision. These results are not used as the final basis for a general improvement claim because the final Tongji results are reported under the stricter development/test subject-disjoint protocol.

\begin{table*}[t]
\centering
\scriptsize
\setlength{\tabcolsep}{3pt}
\caption{Diagnostic Tongji seed42 ablation under the original seen-identity protocol. These results are not the final subject-disjoint evidence.}
\label{tab:ablation_scope}
\resizebox{\textwidth}{!}{
\begin{tabular}{llcccccc}
\toprule
Method & Description & Rank-1 & Rank-5 & Macro-F1 & EER & TAR@1e-2 & TAR@1e-3 \\
\midrule
B1 & ResNet18 + CE + SupCon & 93.39 & 95.47 & 92.43 & 2.30 & 96.23 & 89.75 \\
B4 & ResNet18 + ArcFace & 94.35 & 96.40 & 93.65 & 2.03 & 96.83 & 91.86 \\
B6 & ResNet18 + BNNeck + ArcFace & 96.80 & 97.88 & 96.40 & 1.14 & 98.77 & 96.53 \\
B7 & ResNet18 + BNNeck + ArcFace + light SupCon & 96.97 & 97.98 & 96.60 & 1.13 & 98.77 & 96.50 \\
\bottomrule
\end{tabular}%
}
\end{table*}

The diagnostic ablation shows why B6 was worth evaluating under a stricter protocol. ArcFace and BNNeck improved the reported seen-identity ablation metrics, and light SupCon in B7 did not materially change the result relative to B6. However, this ablation should be read as hyperparameter and architectural evidence, not as proof that B6 is consistently superior.

The subject--disjoint Tongji evaluation in Section~\ref{sec:results} is the decisive evidence for the final paper claim. Under that protocol, B6 does not improve over B1 overall. Therefore, the ablation supports a protocol-sensitivity interpretation rather than a method-dominance claim.

